\documentclass{article}
\usepackage{ctex}
\usepackage{enumitem}
\usepackage{geometry}
\usepackage{xcolor}
\geometry{a4paper, margin=2.5cm}
\begin{document}

\section*{第一章 经济学导论}

\subsection*{一、名词解释}
\begin{enumerate}[label={\arabic*.}]
    \item 亚当·斯密“看不见的手”原理
    \item 稀缺性
\end{enumerate}

\subsection*{二、单项选择题}
\begin{enumerate}[label={\arabic*.}]
    \item 经济学的研究对象是（D）。
    \begin{enumerate}[label={\Alph*.},topsep=0pt,itemsep=0pt]
        \item 如何在股市上赚钱
        \item 证明市场可以配置资源
        \item 选择最公平的收入分配方式
        \item 稀缺资源的配置与利用问题
    \end{enumerate}

    \item 人们在经济资源的配置和利用中要进行选择的根本原因在于（C）。
    \begin{enumerate}[label={\Alph*.},topsep=0pt,itemsep=0pt]
        \item 产品效用的不同
        \item 人们的主观偏好不同
        \item 经济资源的稀缺性
        \item 经济资源用途的不同
    \end{enumerate}

    \item 经济物品可以概括为（D）。
    \begin{enumerate}[label={\Alph*.},topsep=0pt,itemsep=0pt]
        \item 有用的物品
        \item 稀缺的物品
        \item 要用钱购买的物品
        \item 有用且稀缺的物品
    \end{enumerate}

    \item 以下属于自由物品的是（A）。
    \begin{enumerate}[label={\Alph*.},topsep=0pt,itemsep=0pt]
        \item 阳光
        \item 电视机
        \item 汽车
        \item 自来水
    \end{enumerate}

    \item 稀缺性是指（C）。
    \begin{enumerate}[label={\Alph*.},topsep=0pt,itemsep=0pt]
        \item 生产家具的木材太少
        \item 生产粮食的土地太少
        \item 可用的经济资源与人类的欲望相比总是不足
        \item 生产力不发达造成了物品供给不足
    \end{enumerate}

    \item 经济学可以定义为（C）。
    \begin{enumerate}[label={\Alph*.},topsep=0pt,itemsep=0pt]
        \item 研究政府如何对市场机制进行干预的科学
        \item 消费者如何获取收入并进行消费的学说
        \item 研究如何最合理地配置和利用稀缺资源的学科
        \item 研究企业如何取得利润的理论
    \end{enumerate}
\end{enumerate}

\subsection*{三、多项选择题}
\begin{enumerate}[label={\arabic*.}]
    \item 西方经济学教材的特点是（AD）。
    \begin{enumerate}[label={\Alph*.},topsep=0pt,itemsep=0pt]
        \item 偏重理论分析
        \item 纯技术分析
        \item 概括全部经济学流派的观点
        \item 介绍居主流地位的经济学学派的观点
    \end{enumerate}

    \item 一般认为经济学研究的基本问题是（ABCD）。
    \begin{enumerate}[label={\Alph*.},topsep=0pt,itemsep=0pt]
        \item 生产什么，生产多少
        \item 怎样生产
        \item 为谁生产
        \item 何时生产，在哪里生产
    \end{enumerate}

    
    \item 新古典综合派主要由（ABCD）理论综合而成。
    \begin{enumerate}[label={\Alph*.},topsep=0pt,itemsep=0pt]
        \item 1980年马歇尔《经济学原理》
        \item 1933年张伯伦和罗宾逊不完全竞争理论
        \item 1939年希克斯的价值论和一般均衡理论
        \item 1936年凯恩斯《就业、利息和货币通论》
    \end{enumerate}

    \item 资产阶级经济学发展经历的主要阶段有（\quad）。
    \begin{enumerate}[label={\Alph*.},topsep=0pt,itemsep=0pt]
        \item 重商主义经济学
        \item 古典经济学
        \item 庸俗经济学
        \item 庸俗经济学后
    \end{enumerate}

    \item 20世纪70年代边际学派的主要代表人物有（\quad）。
    \begin{enumerate}[label={\Alph*.},topsep=0pt,itemsep=0pt]
        \item 英国学派的杰文斯
        \item 奥地利学派的门格尔
        \item 洛桑学派的瓦尔拉斯
        \item 剑桥学派的马歇尔
    \end{enumerate}

    \item 20世纪30年代西方经济学三次比较重大的修改和补充是（\quad）。
    \begin{enumerate}[label={\Alph*.},topsep=0pt,itemsep=0pt]
        \item 1933年张伯伦和罗宾逊对不完全竞争理论的补充
        \item 1932年萨缪尔森对古典理论的补充
        \item 1936年凯恩斯出版《就业、利息和货币通论》，创立了自成体系的宏观经济学
        \item 1939年希克斯对价值论和一般均衡理论的补充
    \end{enumerate}

    \item 新古典综合派的观点有（ABC）。
    \begin{enumerate}[label={\Alph*.},topsep=0pt,itemsep=0pt]
        \item 倡导宏观调控的市场机制
        \item 倡导私有制为主体的混合所有制
        \item 微观经济理论和宏观经济理论互相补充
        \item 失业和通货膨胀并存
    \end{enumerate}
\end{enumerate}

\subsection*{四、判断题}
    \begin{enumerate}[label={\arabic*.}]
        \item 资源的稀缺性是指世界上的资源终将由于不断用于生产产品和劳务而消耗完。（\quad）
        \item 汽车是可以自由买卖的，所以是自由物品。（\quad）
        \item 1936年凯恩斯发表《就业、利息和货币通论》，标志正式创建了宏观经济理论体系。（\quad）
        \item 萨缪尔森是新古典综合派的杰出代表人物。（\quad）
        \item “交钥匙工程”即把西方工厂丝毫不变地移植到发展中国家就可以带来发展中国家经济快速发展。（\quad）
        \item 资源的有限性和用途的多样性是配置资源的前提条件。（\quad）
        \item “看不见的手”原理是马歇尔在《经济学原理》中首先提出来的。（\quad）
    \end{enumerate}


    \textcolor{red}{自由物品（free goods）是经济学中与“经济物品”相对的概念，指相对与人的欲望而言的供给足够充裕的，无序付出代价即可获得的物品。}
    \textbf {自由物品的核心特点是：\\}
    \textbf{不稀缺：数量足以满足人们在零价格下的需求；\\}
    \textbf{机会成本为零：使用它不需要放弃其他用途；\\}
    \textbf{无需价格配置：不需要通过市场交换或配给来分配。\\}


\end{document}